\documentclass[aspectratio=169,10pt]{beamer}
\usetheme{Boadilla}
\usecolortheme{default}
\definecolor{Ink}{HTML}{14233B}
\definecolor{Blue}{HTML}{2463A6}
\definecolor{Amber}{HTML}{D88922}
\definecolor{Paper}{HTML}{F4F6F9}
\setbeamercolor{normal text}{fg=Ink,bg=white}
\setbeamercolor{structure}{fg=Blue}
\setbeamercolor{frametitle}{fg=Ink,bg=white}
\setbeamercolor{title}{fg=Ink}
\setbeamercolor{subtitle}{fg=Blue}
\setbeamercolor{author}{fg=Ink}
\setbeamercolor{date}{fg=Ink}
\setbeamercolor{codebox}{fg=Ink,bg=Paper}
\setbeamercolor{alerted text}{fg=Amber}
\setbeamertemplate{navigation symbols}{}
\setbeamertemplate{footline}[frame number]
\setbeamersize{text margin left=9mm,text margin right=9mm}
\setbeamertemplate{itemize item}{\color{Amber}\small$\blacktriangleright$}
\setbeamertemplate{itemize subitem}{\color{Blue}\scriptsize$\bullet$}
\usepackage{listings}
\usepackage{amsmath,amssymb}
\lstset{basicstyle=\ttfamily\small,columns=fullflexible,keepspaces=true,
  showstringspaces=false,frame=single,rulecolor=\color{Blue!25},
  backgroundcolor=\color{Paper},xleftmargin=5pt,xrightmargin=5pt,
  aboveskip=7pt,belowskip=7pt}
\newcommand{\meta}[1]{\textcolor{Blue}{\texttt{#1}}}
\newcommand{\kw}[1]{\textcolor{Amber!90!black}{\texttt{#1}}}
\newcommand{\smallnote}[1]{{\small\color{Ink!75}#1}}

\title{Type inference in LAMBDA}
\subtitle{Constraints, annotations, and explicit polymorphism}
\author{CS413 / LAMBDA}
\date{}

\begin{document}

\begin{frame}[plain]
  \titlepage
  \vspace{-5mm}
  \begin{center}\smallnote{From typing constraints to the implementation in \texttt{lambda2.py}}\end{center}
\end{frame}

\begin{frame}[fragile]{Inference as constraint solving}
\begin{columns}[T,onlytextwidth]
\column{.54\textwidth}
Start with unknown types for expressions whose types are not yet known.

Each language construct adds equations between types. Substitution solves the
equations. A conflict means the program is ill typed.

\medskip
The resolved type of the whole expression is the checker result.
\column{.42\textwidth}
\begin{lstlisting}
lam (x) => x + 1

parameter: ?t0
+ requires: ?t0 = int
result: int
function: int -> int
\end{lstlisting}
\end{columns}
\medskip
\smallnote{Omitted types are inferred. Inferred bindings are not generalized.}
\end{frame}

\begin{frame}[fragile]{Types and inference variables}
\begin{columns}[T,onlytextwidth]
\column{.53\textwidth}
\[
\begin{aligned}
T ::=\;& \texttt{int} \mid \texttt{bool} \mid \texttt{unit}\\
 &\mid (T_1,\ldots,T_n) \mid T_1\to T_2\\
 &\mid \texttt{List}[T] \mid \texttt{Optn}[T] \mid ?t_N
\end{aligned}
\]
A type may contain unknowns, for example
\( ?t_0 \to \texttt{List}[?t_1] \).
\column{.43\textwidth}
\begin{itemize}
\item \meta{solver.fresh()} creates an unknown \meta{?tN}.
\item A substitution records what an unknown stands for.
\item \meta{solver.apply(T)} follows substitutions through a type.
\end{itemize}
\end{columns}
\medskip
\begin{minipage}{\textwidth}\smallnote{Concrete types: \texttt{TInt}, \texttt{TBool}, \texttt{TTuple}, \texttt{TFun}, \texttt{TList}, \texttt{TOption}, and \texttt{TMeta}. Declared type variables are represented as \texttt{TPoly} before checking.}\end{minipage}
\end{frame}

\begin{frame}[fragile]{Unification turns uses into type facts}
Two function types unify by matching their arguments and results. Tuple types
unify component by component when their lengths match. Lists and options unify
their element types. A metavariable can be replaced when that does not create
an infinite type.

\begin{columns}[T,onlytextwidth]
\column{.45\textwidth}
\begin{lstlisting}
unify(?t0, int)
  => ?t0 := int
\end{lstlisting}
\column{.52\textwidth}
\begin{lstlisting}
unify(?t1 -> bool, int -> ?t2)
  ?t1 := int
  ?t2 := bool
\end{lstlisting}
\end{columns}
\end{frame}

\begin{frame}[fragile]{A complete inference example}
\begin{center}\Large\ttfamily lam (x) => x + 1\end{center}
\begin{enumerate}
\item Give \texttt{x} a fresh type, \meta{?t0}.
\item The \texttt{+} rule requires both operands to have type \texttt{int}, so \meta{?t0 = int}.
\item The literal \texttt{1} has type \texttt{int}; the body therefore has type \texttt{int}.
\item The lambda type \meta{?t0 -> int} resolves to \texttt{int -> int}.
\end{enumerate}
\end{frame}

\begin{frame}[fragile]{The typing environment}
\begin{columns}[T,onlytextwidth]
\column{.55\textwidth}
\meta{d0exp\_typecheck} checks each syntax tree node. Its environment maps
names to type schemes.
\begin{itemize}
\item A monomorphic binding has no quantified variables.
\item An explicitly polymorphic function quantifies only its declared type parameters.
\item A lambda extends a copied environment with its parameter type.
\end{itemize}
\column{.41\textwidth}
\begin{lstlisting}
inner = dict(env)
inner[param] =
  monomorphic(parameter_type)
\end{lstlisting}
\end{columns}
\medskip
A recursive function binds both its function type and its parameter in the
body. When a declaration gives a function signature, the checker uses it as
context to infer omitted lambda and recursive parameter types.
\end{frame}

\begin{frame}[fragile]{Constraints by construct}
\begin{columns}[T,onlytextwidth]
\column{.48\textwidth}
\textbf{Conditional}
\begin{lstlisting}
if c then a else b
  unify(type(c), bool)
  unify(type(a), type(b))
\end{lstlisting}
\textbf{Application}
\begin{lstlisting}
f(arg)
  unify(type(f), type(arg) -> fresh())
\end{lstlisting}
\column{.48\textwidth}
\textbf{Arithmetic}
\begin{lstlisting}
x + y
  unify(type(x), int)
  unify(type(y), int)
\end{lstlisting}
Container primitives constrain element types. The checker reports
\meta{TypeCheckError} at a conflicting constraint.
\end{columns}
\end{frame}

\begin{frame}[fragile]{Explicit polymorphism, without generalization}
\begin{lstlisting}
identity[T] : T -> T = lam (x) => x
\end{lstlisting}
\begin{itemize}
\item \texttt{T} is declared in the function name's type parameter list.
\item The function needs an explicit signature. The lambda parameter type is inferred from it.
\item Each use instantiates \texttt{T} with a fresh type.
\end{itemize}
\begin{lstlisting}
identity(1)      # T is int at this use
identity(true)   # T is bool at this use
\end{lstlisting}
The same syntax works for inner functions in local declaration blocks. Functions
without explicit type parameters remain monomorphic. Inferred types are never
generalized.
\end{frame}

\begin{frame}[fragile]{Annotations constrain inferred types}
\begin{columns}[T,onlytextwidth]
\column{.47\textwidth}
\begin{lstlisting}
(expression : T)
\end{lstlisting}
\begin{itemize}
\item \texttt{D0Eanno} wraps an expression and its declared type.
\item The checker infers the expression type, then unifies it with \texttt{T}.
\item Declaration signatures also guide inference inside function bodies.
\end{itemize}
\column{.49\textwidth}
\begin{lstlisting}
(lam (x) => x + 1 : int -> int)
# accepted

(1 : bool)
# rejected

id[T] : T -> T = lam (x) => x
# accepted
\end{lstlisting}
\end{columns}
\end{frame}

\begin{frame}[fragile]{What the implementation checks}
\begin{itemize}
\item Applications, arithmetic, comparisons, conditionals, tuples, and projections.
\item Homogeneous list and option operations, with element types inferred from use.
\item Recursive function bodies against their result types.
\item Expression and declaration annotations, including explicit polymorphic signatures.
\item Unbound names, tuple index errors, mismatched types, and infinite types.
\end{itemize}
\begin{lstlisting}
d0exp_parse_and_typecheck(source)
  expression = d0exp_parse(source)
  return expression, d0exp_typecheck(expression)

d0cls_typecheck(d0cls_parse(source))
\end{lstlisting}
\end{frame}

\end{document}
